\documentclass[12pt]{article}

\setlength{\parindent}{0pt}
\setcounter{secnumdepth}{3}
\usepackage[margin=1in]{geometry}

% packages
\usepackage{amsmath, amsfonts, amsthm, amssymb}
\usepackage{graphicx, float}
\usepackage{mathtools}
\usepackage{titlesec}
\usepackage{interval}
\usepackage{hyperref}
\usepackage{siunitx}
\usepackage{titling}
\usepackage{vwcol}
\usepackage{setspace}
\usepackage{empheq}
\usepackage{cancel}
\usepackage{esdiff}
\usepackage{multicol}
\usepackage{mdframed}
\usepackage{esdiff}
\usepackage{tikzsymbols}
\usepackage{multicol}
\usepackage{tikz}
\usepackage{varwidth}
\usepackage{pgfplots}
\usepackage{fancyhdr}
\usepackage{enumitem}
\usepackage{tcolorbox}

\intervalconfig {
	soft open fences
}

% header
\pagestyle{fancy}
\fancyhf{}
\lhead{Ethan Fan}
\rhead{AP Calculus BC}
\cfoot{\thepage}

% section format
\titleformat{\section}
  {\large \bfseries}
  {}
  {0em}
  {}
\titlespacing*{\section}{0pt}{0.5em}{0.3em}

\titleformat{\subsection}[runin]
  {\bfseries}
  {}
  {0em}
  {}
\titlespacing*{\subsection}{0pt}{0pt}{0em}

\titleformat{\subsubsection}[runin]
  {\itshape}
  {(\roman{subsubsection})}
  {0em}
  {}

% commands
\newcommand{\alignedintertext}[1]{%
  \noalign{%
    \vskip\belowdisplayshortskip
    \vtop{\hsize=\linewidth#1\par
    \expandafter}%
    \expandafter\prevdepth\the\prevdepth
  }%
}

\newcommand*{\paren}[1]{\ensuremath\left(#1\right)}
\newcommand*{\limit}[2][x]{\ensuremath{\displaystyle\lim_{#1\to#2}}}
\newcommand*{\Limit}[3][x]{\ensuremath{\displaystyle\lim_{#1\to#2}\left[#3\right]}}
\newcommand*{\deriv}[1][x]{\ensuremath{\dfrac{\mathrm{d}}{\mathrm{d}#1}}}
\newcommand*{\Deriv}[2][x]{\ensuremath{\dfrac{\mathrm{d}}{\mathrm{d}#1}\left[#2\right]}}
\newcommand{\ddx}{\dfrac{d}{dx}}
\newcommand{\dydx}{\dfrac{dy}{dx}}
\newcommand{\dydt}{\dfrac{dy}{dt}}
\newcommand{\dxdt}{\dfrac{dx}{dt}}
\newcommand*{\iinteg}[2][x]{\ensuremath{\displaystyle\int #2\;\mathrm{d}#1}}
\newcommand*{\dinteg}[4][x]{\ensuremath{\displaystyle\int_{#2}^{#3}#4\;\mathrm{d}#1}}
\newcommand*{\abs}[1]{\ensuremath{\left|#1\right|}}
\newcommand*{\eps}{\varepsilon}
\newcommand*{\real}{\ensuremath\mathbb{R}}
\newcommand*{\integer}{\ensuremath\mathbb{Z}}
\newcommand*{\posint}{\ensuremath\mathbb{Z^+}}
\newcommand*{\cbrt}[1]{\ensuremath{\sqrt[3]{#1}}}
\newcommand*{\lheqzero}{\ensuremath{\underset{\text{L'H}}{\overset{\left[\frac00\right]}{=}}}}
\newcommand*{\lheqinfty}{\ensuremath{\underset{\text{L'H}}{\overset{\left[\frac{\infty}{\infty}\right]}{=}}}}
\newcommand{\tor}{\text{ or }}
\newcommand{\floor}[1]{\lfloor #1 \rfloor}

\newcommand{\vem}{\vspace{0.5em}}
\newcommand{\example}[1]{\section{Example #1}}
\newcommand{\problem}[1]{\section{Problem #1}}
\newcommand{\subproblem}[1]{\subsection{(#1)} \,}

\DeclareMathOperator{\dne}{DNE}
\DeclareMathOperator{\sgn}{sgn}

\newcommand{\definition}[1]{\begin{tcolorbox}[colback=red!5!white,colframe=red!75!black,parbox=false] #1 \end{tcolorbox}}

\newcommand{\theorem}[2]{\begin{tcolorbox}[title={#1},colback=blue!5!white,colframe=blue!75!black,parbox=false] #2 \end{tcolorbox}}

\allowdisplaybreaks
% \postdisplaypenalty=100000

% document
\begin{document}

\vspace*{0cm}

% opening
\begin{center}
    {\Large \bfseries Problem Set 9} \\[0.5cm]
    {\large Integration by Substitution} \\[0.35cm]
    {\normalsize September 25, 2026}
\end{center}

% problems

\problem{1}

\subproblem{e}
\begin{gather*}
    u=x^2+1\\
    du=2x\, dx\\
    \int \frac{6x}{x^2+1}\, dx=\boxed{\int \frac{3}{u}\, du}
\end{gather*}

\subproblem{f}
\begin{gather*}
    u=\ln t\\
    du=\frac{1}{t} \, dt\\
    \int \frac{\sqrt{\ln t}}{t}\, dt=\boxed{\int \sqrt{u}\, du}
\end{gather*}

\subproblem{g}
\begin{gather*}
    u=x^2+4x\\
    du=2x+4\, dx\\
    \int (x+2)\tan (x^2+4x)\, dx=\boxed{\frac{1}{2}\int \tan u\, du}
\end{gather*}

\subproblem{h}
\begin{gather*}
    u=\arctan x\\
    du=\frac{1}{x^2+1}\, dx\\
    \int \frac{(\arctan x)^2}{x^2+1}\, dx=\boxed{\int u^2\, du}
\end{gather*}

\problem{2}

\subproblem{e}
\begin{gather*}
    u=2x+1\\
    du=2\, dx\\
    \frac{1}{2}du=dx\\
    \int_0^{13} \frac{1}{\cbrt{(1+2x)^2}}\, dx\\
    =\frac{1}{2}\int_1^{27} \frac{1}{u^\frac{2}{3}}\, du\\
    =\frac{3}{2}(u^\frac{1}{3})\Big|^{27}_1\\
    =\boxed{3}
\end{gather*}

\subproblem{f}
\begin{gather*}
    u=x^2+a^2\\
    du=2x\, dx\\
    \frac{1}{2}du = x\, dx\\
    \int_0^a x\sqrt{x^2+a^2}\, dx\\
    =\frac{1}{2}\int_{a^2}^{2a^2} \sqrt{u}\, du\\
    =\frac{1}{3}(2\sqrt{2}a^3-a^3)\\
    =\boxed{\frac{a^3(2\sqrt{2}-1)}{3}}
\end{gather*}

\subproblem{g}
\begin{gather*}
    u=x-1\\
    du=dx\\
    \int_1^2x\sqrt{x-1}\, dx\\
    =\int_0^1 (u+1)\sqrt{u}\, du\\
    =\int_0^1 u^\frac{3}{2}+u^\frac{1}{2}\, du\\
    = (\frac{2}{5}u^\frac{5}{2}+ \frac{2}{3}u^\frac{3}{2})\Big|^1_0\\
    =\boxed{\frac{16}{15}}
\end{gather*}

\subproblem{h}
\begin{gather*}
    u=\frac{1}{x^2}\\
    du = -\frac{2}{x^3}\, dx\\
    -\frac{1}{2}du=\frac{dx}{x^3}\\
    \int_\frac{1}{2}^1\frac{\cos (x^{-2})}{x^3}\, dx\\
    =\frac{1}{2}\int_1^4 \cos u \, du\\
    =\boxed{\frac{\sin 4-\sin 1}{2}}
\end{gather*}

\problem{5}

\subproblem{a}
\begin{gather*}
    f(x)=2x\cos x^2\\
    \cos x^2=0 \implies x_k^2=\frac{\pi}{2}+k\pi\\
    x_k=\sqrt{\frac{(2k+1)\pi}{2}}\\
    \lim_{k\to \infty} x_{k+1}-x_k=0
\end{gather*}

\subproblem{b}
\begin{gather*}
    u=x^2\\
    du=2x \, dx\\
    \int 2x\cos x^2\, dx=\int \cos u\, du\\
    \int_0^{\sqrt{\frac{\pi}{2}}} 2x\cos x^2\, dx=\int_0^\frac{\pi}{2} \cos u\, du = 1\\
    \int_{x_k}^{x_{k+1}} |2x\cos x^2|\, dx = \int_{\frac{(2k+1)\pi}{2}}^{\frac{(2k+3)\pi}{2}} |\cos u|\, du = \boxed{2}
\end{gather*}
The area of every full lobe is 2.\vem

\subproblem{c}
The lobes of the original function are all stretched to the lobes of the standard function $\cos u$, which is always 2.\vem

\subproblem{d}
\[
\int_0^{x_{2k}} 2x\cos x^2\, dx= \boxed{1}\\
\]

\problem{6}

\subproblem{a}
\begin{gather*}
    u=-x\\
    du=-dx\\
    \int_{-a}^0f(x)\, dx=-\int_{-a}^0f(-u)\, du=\int_{0}^af(-u)\, du
\end{gather*}

\subproblem{b}
\begin{gather*}
    \int_{-a}^af(x)\, dx=\int_{0}^af(x)\, dx+\int_{0}^af(-x)\, dx=2\int_{0}^af(x)\, dx\\
    \int_{-a}^af(x)\, dx=\int_{0}^af(x)\, dx+\int_{0}^af(-x)\, dx=0
\end{gather*}

\subproblem{c}
\begin{gather*}
    \int_{-2}^2x^3\cos x^2\, dx=0\\
     \int_{-\pi/2}^{\pi/2} (x^2 \sin x + \cos x) \, dx=0+     \int_{-\pi/2}^{\pi/2} \cos x\, dx=2
\end{gather*}

\subproblem{d}
No, as the original function is not continuous on these intervals.

\problem{7}

\subproblem{a}
\begin{gather*}
    u=\sin x\\
    du =\cos x\, dx\\
    \int \sin x \cos x \, dx = \int u\, dx= \boxed{\frac{\sin^2 x}{2}+C}
\end{gather*}

\subproblem{b}
\begin{gather*}
    u=\cos x\\
    du =-\sin  x\, dx\\
    \int \sin x \cos x \, dx = -\int u\, dx=\boxed{-\frac{\cos^2x}{2}+C}
\end{gather*}

\subproblem{c}
\begin{gather*}
    \int \sin x \cos x \, dx= \frac{1}{2}\int \sin 2x\, dx \\
    u=2x\\
    du=2dx\\
    \frac{1}{2}\int \sin 2x\, dx =\frac{1}{4}\int \sin u \, du=\boxed{-\frac{\cos 2x}{4}+C}
\end{gather*}

\subproblem{d}
As they merely differ by a constant, all three are correct anti-derivatives.

\problem{9}

\subproblem{e}
No possible $u$-sub exists as the result of integration changes.\vem

\subproblem{f}
\begin{gather*}
    u=\sin x\\
    du=\cos x\, dx\\
    \int_0^\pi e^{\sin x}\cos x\, dx =\int_0^0e^u\, du=0
\end{gather*}
The integral evaluates to zero as the function is symmetric about $x=\dfrac{\pi}{2}$.

\problem{10}

\subproblem{a}
\begin{gather*}
    T=\frac{1}{60}\\
    V_{avg} = 60\int_0^{\frac{1}{60}}170\sin (120\pi t)=0 
\end{gather*}

\subproblem{b}
\begin{gather*}
    V_{avg} = 170\cdot 120\int_0^{\frac{1}{120}}\sin (120\pi t)=170\cdot 120(-\frac{\cos 120\pi t}{120\pi})\Big|^{\frac{1}{120}}_0=\boxed{\frac{340}{\pi}}\\
    \frac{340}{\pi} \approx \boxed{108.23}
\end{gather*}

\subproblem{c}
\begin{gather*}
    v^2(t)=170^2\sin ^2(120\pi t)=\frac{170^2}{2}(1-\cos 240\pi t)\\
    v^2_{avg}=60\cdot \frac{170^2}{2}\int_0^\frac{1}{60} 1-\cos 240\pi t \, dt=30\cdot 170^2(t-\frac{\sin 240\pi t}{240})\Big|_0^\frac{1}{60}=14450\\
    V_{rms}=\sqrt{14450}=\boxed{85\sqrt{2}\approx 120.2}
\end{gather*}

\subproblem{d}
\begin{gather*}
    T=\frac{2\pi}{\omega }\\
    V^2_{avg}=\frac{2\pi}{\omega} \int_0^\frac{\omega}{2\pi } A^2\sin ^2\omega t\, dt\\
    =\frac{A^2\omega}{4\pi }\int_0^\frac{2\pi}{\omega } 1-\cos 2\omega t\, dt\\
    =\frac{A^2\omega}{4\pi } [t-\frac{\sin 2\omega t}{2\omega }]\Big|^\frac{2\pi}{\omega}_0\\
    =\frac{A^2\omega}{4\pi } \cdot \frac{2\pi}{\omega}= \frac{A^2}{2}\\
    \boxed{V_{rms}=\frac{A}{\sqrt{2}}}
\end{gather*}

\subproblem{e}
The label "120 V" refers to the answer to (c). By utilizing the $V_{rms}$, engineers may calculate values of the average power over certain intervals. 

\problem{11}

\subproblem{a}
\begin{gather*}
    u=a-x\\
    du=-dx\\
    \int_0^af(x)\, dx=\int_a^0f(a-u)\, -du=\int_0^af(a-u)\, du=\int_0^af(a-x)\,d x
\end{gather*}

\subproblem{b}
\begin{gather*}
    I=\int_0^\frac{\pi}{2}\frac{\sin^{2026}x}{\sin^{2026} x+\cos^{2026}x}dx\\=\int_0^\frac{\pi}{2}\frac{\sin^{2026}(\frac{\pi}{2}-x)}{\sin^{2026} (\frac{\pi}{2}-x)+\cos^{2026}(\frac{\pi}{2}-x)}dx=\int_0^\frac{\pi}{2}\frac{\cos^{2026}x}{\sin^{2026} x+\cos^{2026}x}dx\\
    2I=\int_0^\frac{\pi}{2}\frac{\sin^{2026} x+\cos^{2026}x}{\sin^{2026} x+\cos^{2026}x}dx=\int_0^\frac{\pi}{2}1\, dx=\frac{\pi}{2}\\
    \boxed{I=\frac{\pi}{4}}
\end{gather*}

\subproblem{c}
\begin{gather*}
    \int_0^\pi x\, g(\sin x)\, dx= \int_0^\pi (\pi -x)\, g(\sin \pi -x)\, dx\\=\pi \int_0^\pi g(\sin x)\,d x-\int_0^\pi x\, g(\sin x)\, dx\\
    2\int_0^\pi x\, g(\sin x)\, dx=\pi \int_0^\pi \, g(\sin x)\, dx\\
    \int_0^\pi x\, g(\sin x)\, dx=\frac{\pi}{2} \int_0^\pi \, g(\sin x)\, dx\\
    \int_0^\pi \frac{x\sin x}{1+\cos^2x}\, dx=\frac{\pi}{2} \int_0^\pi \frac{\sin x}{1+\cos^2 x}\, dx\\
    u=\cos x\\
    du=-\sin x \, dx\\
    \frac{\pi}{2} \int_0^\pi \frac{\sin x}{1+\cos^2 x}\, dx=\frac{\pi}{2} \int_{-1}^1 \frac{1}{1+u^2}\, du=\frac{\pi}{2} \arctan u \Big|^1_{-1}=\boxed{\frac{\pi^2}{4}}
\end{gather*}

\subproblem{d}
\begin{gather*}
    \tan u=x\\
    \sec^2 u\, du=dx\\
    \int_0^1\frac{\ln(1+x)}{1+x^2}\,dx =\int_0^\frac{\pi}{4} \frac{\ln (1+\tan u)}{1+\tan^2u}\cdot \sec^2u\, du=\int_0^\frac{\pi}{4}\ln (1+\tan u)\, du\\
    \int_0^\frac{\pi}{4}\ln (1+\tan u)\, du=\int_0^\frac{\pi}{4}\ln (1+\tan (\frac{\pi}{4}-u))\, du = \int_0^\frac{\pi}{4}\ln (\frac{2}{1+\tan u})\, du\\
    \int_0^\frac{\pi}{4}\ln (\frac{2}{1+\tan u})\, du=\int_0^\frac{\pi}{4}\ln 2 \, du-\int_0^\frac{\pi}{4} \ln (1+\tan u)\, du=\int_0^\frac{\pi}{4} \ln (1+\tan u)\, du\\
    \int_0^\frac{\pi}{4} \ln (1+\tan u)\, du= \frac{1}{2}\cdot \frac{\pi}{4}\cdot \ln 2 = \boxed{\frac{\pi}{8}\ln 2}
\end{gather*}















\end{document}

